\documentclass[conference,10pt]{IEEEtran}
\usepackage[T1]{fontenc}
\usepackage[utf8]{inputenc}
\usepackage[spanish,es-nodecimaldot]{babel}
\usepackage{amsmath}
\usepackage{circuitikz}
\usepackage{xcolor}
\newcommand{\pendiente}[1]{\textcolor{red}{\textit{[Completar: #1]}}}
\newcommand{\paralelo}{\mathbin{\|}}

\title{Preinforme: Circuitos Resistivos -- Parte 1\\Práctica de Laboratorio No. 2}
\author{\IEEEauthorblockN{\pendiente{Nombre completo 1}, \pendiente{Nombre completo 2}}
\IEEEauthorblockA{Ingeniería Electrónica y de Telecomunicaciones\\
Circuitos Eléctricos I -- Universidad de Antioquia\\2026-2\\
\pendiente{Correos institucionales}}}

\begin{document}
\maketitle
\begin{abstract}
Este preinforme presenta la preparación para la Práctica de Laboratorio No. 2,
\emph{Circuitos Resistivos -- Parte 1}. Se documentan los componentes
seleccionados, la consulta sobre transformaciones delta--estrella y
estrella--delta, y el cálculo del resistor equivalente y la corriente de una
fuente de $5\,\mathrm{V}$ en tres circuitos resistivos.
\end{abstract}
\begin{IEEEkeywords}
Ley de Ohm, resistores, serie, paralelo, delta--estrella, estrella--delta.
\end{IEEEkeywords}

\section{Objetivos}
\begin{itemize}
  \item Realizar el montaje de circuitos resistivos en serie, paralelo y mixtos.
  \item Verificar la Ley de Ohm y la simplificación de circuitos mediante
  conexiones serie/paralelo y transformaciones delta--estrella.
\end{itemize}

\section{Componentes y valores seleccionados}
Se llevará una protoboard, cables de conexión, multímetro, fuente de voltaje y
cinco resistores comerciales diferentes entre $330\,\Omega$ y
$4.7\,\mathrm{k}\Omega$. La fuente se ajustará a $5\,\mathrm{V}$.
\begin{table}[!t]
\caption{Resistores seleccionados para el laboratorio}
\label{tab:resistores}
\centering
\begin{tabular}{ccc}
\hline
Resistor & Valor nominal & Valor medido \\
\hline
$R_a$ & \pendiente{} & \pendiente{} \\
$R_b$ & \pendiente{} & \pendiente{} \\
$R_c$ & \pendiente{} & \pendiente{} \\
$R_d$ & \pendiente{} & \pendiente{} \\
$R_e$ & \pendiente{} & \pendiente{} \\
\hline
\end{tabular}
\end{table}

\section{Consulta: transformaciones delta--estrella y estrella--delta}
Las transformaciones delta--estrella y estrella--delta sustituyen una red de
tres resistores por otra red equivalente vista desde los mismos tres
terminales. Se emplean cuando un circuito no puede reducirse solo mediante
asociaciones en serie o paralelo. En la transformación delta--estrella, cada
resistor de la estrella es el producto de los dos resistores delta que llegan
al mismo terminal, dividido entre la suma de los tres resistores delta. En la
transformación inversa, cada resistor delta es la suma de los productos de dos
resistores estrella divididos por el resistor estrella opuesto. Ambas redes
preservan la resistencia equivalente entre cualquier par de terminales
\cite{sadiku}.

\subsection{Ecuaciones de transformación}
Para una delta con resistores $R_{ab}$, $R_{bc}$ y $R_{ca}$, defina
\begin{equation}
R_\Delta=R_{ab}+R_{bc}+R_{ca}.
\end{equation}
La estrella equivalente se obtiene con
\begin{align}
R_a&=\frac{R_{ab}R_{ca}}{R_\Delta}, &
R_b&=\frac{R_{ab}R_{bc}}{R_\Delta}, &
R_c&=\frac{R_{bc}R_{ca}}{R_\Delta}.
\end{align}
Para una estrella con resistores $R_a$, $R_b$ y $R_c$, las resistencias delta
equivalentes son
\begin{align}
R_{ab}&=R_a+R_b+\frac{R_aR_b}{R_c},\\
R_{bc}&=R_b+R_c+\frac{R_bR_c}{R_a},\\
R_{ca}&=R_c+R_a+\frac{R_cR_a}{R_b}.
\end{align}

\section{Cálculos previos}
En los tres casos se toma $v_s=5\,\mathrm{V}$. Reemplace los campos pendientes
por los valores nominales escogidos y muestre el desarrollo algebraico antes
de registrar cada resultado.

\subsection{Circuito en paralelo}
\begin{figure}[!t]
\centering
\begin{circuitikz}[american voltages]
\draw (0,0) to[battery1,l=$v_s$] (0,2.6) -- (5.4,2.6) -- (5.4,0) -- (0,0);
\draw (1.4,2.6) to[R,l=$R_a$] (1.4,0);
\draw (2.7,2.6) to[R,l=$R_b$] (2.7,0);
\draw (4.0,2.6) to[R,l=$R_c$] (4.0,0);
\draw[-latex] (0.35,2.3) -- (1.0,2.3) node[midway,above] {$i$};
\end{circuitikz}
\caption{Circuito resistivo en paralelo (Figura 2 de la guía).}
\label{fig:paralelo}
\end{figure}
\begin{align}
R_{\mathrm{eq}1}&=\left(\frac{1}{R_a}+\frac{1}{R_b}+\frac{1}{R_c}\right)^{-1},&
i&=\frac{v_s}{R_{\mathrm{eq}1}}.
\end{align}
\noindent\textbf{Resultado:}\quad
$R_{\mathrm{eq}1}=\pendiente{}$,\quad $i=\pendiente{}$.

\subsection{Circuito en serie}
\begin{figure}[!t]
\centering
\begin{circuitikz}[american voltages]
\draw (0,0) to[battery1,l=$v_s$] (0,2.2)
to[R,l=$R_a$] (2.0,2.2) to[R,l=$R_b$] (4.0,2.2)
to[R,l=$R_c$] (6.0,2.2) -- (6.0,0) -- (0,0);
\draw[-latex] (0.35,1.9) -- (0.95,1.9) node[midway,above] {$i$};
\end{circuitikz}
\caption{Circuito resistivo en serie (Figura 3 de la guía).}
\label{fig:serie}
\end{figure}
\begin{align}
R_{\mathrm{eq}2}&=R_a+R_b+R_c,& i&=\frac{v_s}{R_{\mathrm{eq}2}}.
\end{align}
\noindent\textbf{Resultado:}\quad
$R_{\mathrm{eq}2}=\pendiente{}$,\quad $i=\pendiente{}$.

\subsection{Circuito que requiere transformación delta--estrella}
\begin{figure}[!t]
\centering
\begin{circuitikz}[american voltages]
\draw (0,0) to[battery1,l=$v_s$] (0,4.1) -- (4.4,4.1) -- (4.4,0) -- (0,0);
\draw (1.5,4.1) to[R,l=$R_a$] (1.5,2.05) to[R,l=$R_d$] (1.5,0);
\draw (3.2,4.1) to[R,l=$R_b$] (3.2,2.05) to[R,l=$R_e$] (3.2,0);
\draw (1.5,2.05) to[R,l=$R_c$] (3.2,2.05);
\draw[-latex] (0.35,3.75) -- (0.95,3.75) node[midway,above] {$i$};
\end{circuitikz}
\caption{Puente resistivo que requiere una transformación (Figura 4 de la guía).}
\label{fig:deltaestrella}
\end{figure}
La delta formada por $R_a$, $R_b$ y $R_c$ se transforma en una estrella:
\begin{align}
R_T&=\frac{R_aR_b}{R_a+R_b+R_c},\\
R_L&=\frac{R_aR_c}{R_a+R_b+R_c},\\
R_R&=\frac{R_bR_c}{R_a+R_b+R_c}.
\end{align}
Luego,
\begin{align}
R_{\mathrm{eq}3}&=R_T+\left(R_L+R_d\right)\paralelo\left(R_R+R_e\right),&
i&=\frac{v_s}{R_{\mathrm{eq}3}}.
\end{align}
\noindent\textbf{Desarrollo numérico:}\quad
\pendiente{Sustituir $R_a$ a $R_e$ y mostrar las operaciones.}

\medskip
\noindent\textbf{Resultado:}\quad
$R_{\mathrm{eq}3}=\pendiente{}$,\quad $i=\pendiente{}$.

\section{Conclusión preliminar}
Se espera contrastar los resultados calculados con los valores de simulación y
las mediciones del montaje. Las diferencias se analizarán considerando la
tolerancia de los resistores, la precisión del multímetro y las resistencias de
contacto de la protoboard.

\begin{thebibliography}{1}
\bibitem{guia}
Universidad de Antioquia, Práctica de Laboratorio No. 2: Circuitos
Resistivos -- Parte 1, guía de laboratorio, Circuitos Eléctricos I, 2026-2.
\bibitem{sadiku}
C. K. Alexander y M. N. O. Sadiku, \emph{Fundamentals of Electric Circuits},
5th ed. New York, NY, USA: McGraw-Hill, 2013.
\end{thebibliography}
\end{document}
